% GENERATED FILE. Edit canonical lessons, then run npm run generate:books.
% canonical-source-hash: fnv1a64:da5715028e465f7d
% canonical-lessons: HI-C318-dhamaka, HI-C318-bam, HI-C318-gujna, HI-C318-tharrana, HI-C318-sadma

\chapter{A Blast}
\label{ch:hi-a-blast}
\hlchaptermodality{318}

\begin{chapteropening}
\textbf{By the end of this chapter:} \emph{I can describe a sudden loud event.}

Complete the last lesson of chapter 318: I can describe a sudden loud event.
\end{chapteropening}

\section[\texorpdfstring{\hi{धमाका} (dhamākā)}{धमाका (dhamākā)}]{\hi{धमाका} (dhamākā) — a bang; a blast}

\label{lesson:HI-C318-dhamaka}



\begin{quote}
Before the new one: say the Hindi for to be startled, then the Hindi for astonished.
\end{quote}

\subsection*{You'll want to know: \hi{धमाका}}
\textbf{\hi{धमाका}} — \emph{dhamākā} — \textquotedblleft{}a bang; a blast\textquotedblright{}.

\begin{quote}
\hi{रात} \hi{को} \hi{एक} \hi{ज़ोर} \hi{का} \hi{धमाका} \hi{हुआ।}
\end{quote}

\begin{quote}
\emph{rāt ko ek zor kā dhamākā huā.}
\end{quote}

\begin{quote}
There was a loud bang in the night.
\end{quote}

Masculine.

\begin{cousinweb}
\textbf{\hi{धमाका}} (\emph{dhamākā}) is probably imitative, built on \emph{dham}, the sound of a heavy thud. CDIAL reconstructs \emph{*dhamm}, \textquotedblleft{}noise\textquotedblright{} (CDIAL 6735), and lists Hindi \emph{dhamādham}, noisily, under it, though not this word itself. \textbf{\hi{धमाका} \hi{होना}} is for there to be a bang.
\end{cousinweb}

\begin{grammarlens}[title={What you've built}]
One more word for telling what happened.
\end{grammarlens}

\subsection*{Your turn}
\begin{itemize}
\raggedright
  \item \emph{Say it:} \emph{dhamākā}
  \item \emph{Say it:} \emph{dhamākā}, once more
  \item \emph{Say it:} \emph{rāt ko ek zor kā dhamākā huā.} — There was a loud bang in the night.
  \item \emph{Recall:} say the Hindi for to be startled, then the Hindi for astonished, then say \emph{dhamākā} again
\end{itemize}

\begin{tcolorbox}[breakable,colback=teal!4,colframe=teal!35!black,title={Before you move on}]
What does \hi{धमाका} mean? (\textquotedblleft{}A bang; a blast\textquotedblright{}.) Say it once more.
\end{tcolorbox}
\section[\texorpdfstring{\hi{बम} (bam)}{बम (bam)}]{\hi{बम} (bam) — a bomb}

\label{lesson:HI-C318-bam}



\begin{quote}
Before the new one: say the Hindi for astonished, then the Hindi for a bang.
\end{quote}

\subsection*{You'll want to know: \hi{बम}}
\textbf{\hi{बम}} — \emph{bam} — \textquotedblleft{}a bomb\textquotedblright{}.

\begin{quote}
\hi{पुलिस} \hi{को} \hi{स्टेशन} \hi{के} \hi{पास} \hi{एक} \hi{बम} \hi{मिला।}
\end{quote}

\begin{quote}
\emph{pulis ko sṭeśan ke pās ek bam milā.}
\end{quote}

\begin{quote}
The police found a bomb near the station.
\end{quote}

Masculine.

\begin{cousinweb}
\textbf{\hi{बम}} (\emph{bam}) is borrowed from English \emph{bomb}, shortened to one Hindi syllable and masculine.
\end{cousinweb}

\begin{grammarlens}[title={What you've built}]
One more word for telling what happened.
\end{grammarlens}

\subsection*{Your turn}
\begin{itemize}
\raggedright
  \item \emph{Say it:} \emph{bam}
  \item \emph{Say it:} \emph{bam}, once more
  \item \emph{Say it:} \emph{pulis ko sṭeśan ke pās ek bam milā.} — The police found a bomb near the station.
  \item \emph{Recall:} say the Hindi for astonished, then the Hindi for a bang, then say \emph{bam} again
\end{itemize}

\begin{tcolorbox}[breakable,colback=teal!4,colframe=teal!35!black,title={Before you move on}]
What does \hi{बम} mean? (\textquotedblleft{}A bomb\textquotedblright{}.) Say it once more.
\end{tcolorbox}
\section[\texorpdfstring{\hi{गूँजना} (gū̃jnā)}{गूँजना (gū̃jnā)}]{\hi{गूँजना} (gū̃jnā) — to echo; to resound}

\label{lesson:HI-C318-gujna}



\begin{quote}
Before the new one: say the Hindi for a bang, then the Hindi for a bomb.
\end{quote}

\subsection*{You'll want to know: \hi{गूँजना}}
\textbf{\hi{गूँजना}} — \emph{gū̃jnā} — \textquotedblleft{}to echo; to resound\textquotedblright{}.

\begin{quote}
\hi{धमाका} \hi{पूरे} \hi{शहर} \hi{में} \hi{गूँजा।}
\end{quote}

\begin{quote}
\emph{dhamākā pūre śahar meṁ gū̃jā.}
\end{quote}

\begin{quote}
The blast echoed through the whole city.
\end{quote}

Past (it …): \hi{वह} \hi{घंटा} \hi{गूँजा}, \emph{vah ghaṇṭā gū̃jā}.

\begin{cousinweb}
\textbf{\hi{गूँजना}} (\emph{gū̃jnā}) goes back to Sanskrit \emph{guñjati}, \textquotedblleft{}hums, buzzes\textquotedblright{} (CDIAL 4175). A sound \textbf{\hi{गूँजता} \hi{है}} when it fills a place and comes back.
\end{cousinweb}

\begin{grammarlens}[title={What you've built}]
One more word for telling what happened.
\end{grammarlens}

\subsection*{Your turn}
\begin{itemize}
\raggedright
  \item \emph{Say it:} \emph{gū̃jnā}
  \item \emph{Say it:} \emph{gū̃jnā}, once more
  \item \emph{Say it:} \emph{dhamākā pūre śahar meṁ gū̃jā.} — The blast echoed through the whole city.
  \item \emph{Recall:} say the Hindi for a bang, then the Hindi for a bomb, then say \emph{gū̃jnā} again
\end{itemize}

\begin{tcolorbox}[breakable,colback=teal!4,colframe=teal!35!black,title={Before you move on}]
What does \hi{गूँजना} mean? (\textquotedblleft{}To echo; to resound\textquotedblright{}.) Say it once more.
\end{tcolorbox}
\section[\texorpdfstring{\hi{थर्राना} (tharrānā)}{थर्राना (tharrānā)}]{\hi{थर्राना} (tharrānā) — to quake; to shudder}

\label{lesson:HI-C318-tharrana}



\begin{quote}
Before the new one: say the Hindi for a bomb, then the Hindi for to echo.
\end{quote}

\subsection*{You'll want to know: \hi{थर्राना}}
\textbf{\hi{थर्राना}} — \emph{tharrānā} — \textquotedblleft{}to quake; to shudder\textquotedblright{}.

\begin{quote}
\hi{धमाके} \hi{से} \hi{खिड़कियाँ} \hi{थर्रा} \hi{गईं।}
\end{quote}

\begin{quote}
\emph{dhamāke se khiṛkiyā̃ tharrā gaīṁ.}
\end{quote}

\begin{quote}
The windows shook with the blast.
\end{quote}

Past (it …): \hi{वह} \hi{दीवार} \hi{थर्रा} \hi{गई}, \emph{vah dīvār tharrā gaī}.

\begin{cousinweb}
\textbf{\hi{थर्राना}} (\emph{tharrānā}) is inherited from \emph{*thar}, a form CDIAL reconstructs for trembling (CDIAL 6092). It is a hard shaking, of a building or of a person in fear.
\end{cousinweb}

\begin{grammarlens}[title={What you've built}]
One more word for telling what happened.
\end{grammarlens}

\subsection*{Your turn}
\begin{itemize}
\raggedright
  \item \emph{Say it:} \emph{tharrānā}
  \item \emph{Say it:} \emph{tharrānā}, once more
  \item \emph{Say it:} \emph{dhamāke se khiṛkiyā̃ tharrā gaīṁ.} — The windows shook with the blast.
  \item \emph{Recall:} say the Hindi for a bomb, then the Hindi for to echo, then say \emph{tharrānā} again
\end{itemize}

\begin{tcolorbox}[breakable,colback=teal!4,colframe=teal!35!black,title={Before you move on}]
What does \hi{थर्राना} mean? (\textquotedblleft{}To quake; to shudder\textquotedblright{}.) Say it once more.
\end{tcolorbox}
\section[\texorpdfstring{\hi{सदमा} (sadmā)}{सदमा (sadmā)}]{\hi{सदमा} (sadmā) — a shock (of grief or fright)}

\label{lesson:HI-C318-sadma}



\begin{quote}
Before the new one: say the Hindi for to echo, then the Hindi for to quake.
\end{quote}

\subsection*{You'll want to know: \hi{सदमा}}
\textbf{\hi{सदमा}} — \emph{sadmā} — \textquotedblleft{}a shock (of grief or fright)\textquotedblright{}.

\begin{quote}
\hi{धमाके} \hi{के} \hi{बाद} \hi{मीरा} \hi{को} \hi{गहरा} \hi{सदमा} \hi{लगा।}
\end{quote}

\begin{quote}
\emph{dhamāke ke bād Mīrā ko gahrā sadmā lagā.}
\end{quote}

\begin{quote}
After the blast Mira was in deep shock.
\end{quote}

Masculine.

\begin{cousinweb}
\textbf{\hi{सदमा}} (\emph{sadmā}) is usually traced, through Persian, to Arabic \emph{ṣadma}, \textquotedblleft{}a blow, a collision\textquotedblright{}: a blow to the mind. A shock \textbf{\hi{लगता} \hi{है}}: \textbf{\hi{मीरा} \hi{को} \hi{सदमा} \hi{लगा}}.
\end{cousinweb}

\begin{grammarlens}[title={What you've built}]
One more word for telling what happened.
\end{grammarlens}

\subsection*{Your turn}
\begin{itemize}
\raggedright
  \item \emph{Say it:} \emph{sadmā}
  \item \emph{Say it:} \emph{sadmā}, once more
  \item \emph{Say it:} \emph{dhamāke ke bād Mīrā ko gahrā sadmā lagā.} — After the blast Mira was in deep shock.
  \item \emph{Recall:} say the Hindi for to echo, then the Hindi for to quake, then say \emph{sadmā} again
\end{itemize}

\begin{tcolorbox}[breakable,colback=teal!4,colframe=teal!35!black,title={Before you move on}]
What does \hi{सदमा} mean? (\textquotedblleft{}A shock (of grief or fright)\textquotedblright{}.) Say it once more.
\end{tcolorbox}
